% ============================================================
% ECEN 662 Project Report
% Compile: pdflatex proposal && bibtex proposal && pdflatex proposal && pdflatex proposal
% Requires neurips_2026.sty in the same directory
% ============================================================

\documentclass{article}
\usepackage[final]{neurips_2026}

\usepackage[utf8]{inputenc}
\usepackage[T1]{fontenc}
\usepackage{amsmath, amssymb, amsfonts}
\usepackage{algorithm}
\usepackage{algpseudocode}
\usepackage{booktabs}
\usepackage{graphicx}
\usepackage{natbib}
\usepackage{hyperref}
\usepackage{pgfplots}
\usepackage{xcolor}


\makeatletter
\def\@noticestring{}
\makeatother

% Method palette (matches the Python script's METHOD_COLORS).
\definecolor{baseColor}{HTML}{1F77B4}
\definecolor{vemColor}{HTML}{D62728}
\definecolor{ecmColor}{HTML}{2CA02C}
\definecolor{gemColor}{HTML}{9467BD}

% Shared pgfplots style for all EM panels.
\pgfplotsset{
  emaxis/.style={
    width=0.48\linewidth,
    height=0.36\linewidth,
    grid=both,
    grid style={dashed, gray!30},
    tick label style={font=\footnotesize},
    label style={font=\small},
    title style={font=\small, yshift=-1ex},
    legend cell align={left},
    legend style={font=\footnotesize, fill=white, fill opacity=0.85,
                  draw=none, inner sep=2pt},
  },
}

% Band+median line driven by columns <prefix>_p25, <prefix>_med, <prefix>_p75
% in the data file passed as the first argument.
% #1 file, #2 color, #3 prefix, #4 x-column key, #5 legend label
\newcommand{\embandline}[5]{%
  \addplot[name path=#3lo, draw=none, forget plot]
    table[x=#4, y=#3_p25, col sep=comma] {#1};
  \addplot[name path=#3hi, draw=none, forget plot]
    table[x=#4, y=#3_p75, col sep=comma] {#1};
  \addplot[#2, opacity=0.18, draw=none, forget plot]
    fill between[of=#3lo and #3hi];
  \addplot[#2, very thick]
    table[x=#4, y=#3_med, col sep=comma] {#1};
  \addlegendentry{#5}%
}

\newcommand{\E}{\mathbb{E}}
\newcommand{\KL}{\mathrm{KL}}

\title{A Comparative Study of EM Variants}

\author{%
  Brody Jordan \\
  Department of Electrical \& Computer Engineering \\
  Texas A\&M University \\
  \texttt{blj@tamu.edu} \\
}

\begin{document}
\maketitle
\pgfplotsset{compat=1.18}
\usepgfplotslibrary{groupplots}
\usepgfplotslibrary{fillbetween}

% ============================================================
\section{Introduction}
% ============================================================

\subsection{Topic Selection}
While reviewing the lecture notes for this course, I found the EM algorithm section particularly interesting. Notably,
in the later slides, a few different variants of the EM algorithm were briefly mentioned, but not explored in depth. I
decided to implement and compare these variants on a shared testbed to better understand their operational
differences. So while there is
some experimentation, I think it is closer to a reading project.

\subsection{Problem Overview}
The central challenge in latent variable models is that the marginal
likelihood requires integrating out the unobserved variable $U$:
\begin{equation}
  f_\theta(x) = \int f_\theta(x, u)\, du.
  \label{eq:marginal}
\end{equation}
This integral appears inside the log in the maximum likelihood
estimation (MLE) objective, making direct maximization infeasible.
EM resolves this via successive approximation, but
compute requirements can still be a bottleneck for more complex models.
% This project
% studies four strategies for handling that bottleneck on a shared Hidden
% Markov Model (HMM) testbed.

% ============================================================
\section{Background}
% ============================================================

\subsection{Maximum Likelihood Estimation with Latent Variables}

Given observations $X$ and latent variables $U$, the MLE problem is:
\begin{equation}
  \hat\theta(x) \in \arg\max_{\theta \in \Theta}\; \log f_\theta(x),
  \qquad f_\theta(x) = \int f_\theta(x, u)\, du.
\end{equation}
The marginalization inside the log is the core difficulty.

\subsection{The Evidence Lower Bound (ELBO)}

For any auxiliary distribution $q(u)$, the log-likelihood decomposes as:
\begin{equation}
  \log f_\theta(x)
  \;=\;
  \underbrace{
    \E_q\!\left[\log f_\theta(x, u)\right] - \E_q\!\left[\log q(u)\right]
  }_{\displaystyle F(q,\, \theta)\ (\text{ELBO})}
  \;+\;
  \KL\!\left(q \;\|\; f_\theta(u \mid x)\right)
  \;\geq\; F(q,\, \theta).
  \label{eq:elbo}
\end{equation}
Since $\KL \geq 0$, the ELBO is always a lower bound on the log-likelihood.
Setting $q = f_\theta(u \mid x)$ collapses the KL term to zero and makes the
bound tight.

\subsection{Computational Complexity}

For a Gaussian Mixture Model with $K$ components and $n$ observations, the
evidence requires summing over all $K^n$ discrete assignment configurations:
\begin{equation}
  m(x) = \sum_c p(c) \int f(\mu) \prod_{i=1}^n f(x_i \mid \mu_{c_i})\, d\mu.
\end{equation}
This scales exponentially in $n$, making exact inference strictly intractable.


% ============================================================
\section{Methods}
% ============================================================

All four methods alternate between an E-step (computing or approximating
the posterior over latent states) and an M-step (updating $\theta$). They
differ in how the E-step and M-step are handled.

\subsection{Base EM}

Beginning with the base EM algorithm, we define:
\begin{align}
  \alpha_t(k) & = f_\theta(x_1, \ldots, x_t,\, Z_t = k),       \\
  \beta_t(k)  & = f_\theta(x_{t+1}, \ldots, x_T \mid Z_t = k).
\end{align}
The posterior state marginals and pairwise marginals are
\begin{align}
  \gamma_t(k) & = \frac{\alpha_t(k)\,\beta_t(k)}{\sum_j \alpha_t(j)\,\beta_t(j)},
  \label{eq:gamma}                                                                             \\
  \xi_t(j, k) & \propto \alpha_t(j)\,A_{jk}\,f_\theta(x_{t+1} \mid Z_{t+1}=k)\,\beta_{t+1}(k).
  \label{eq:xi}
\end{align}
With $q^{(t)}(z) = f_{\theta^{(t-1)}}(z \mid x)$ the $\KL$ term in
Eq.~\eqref{eq:elbo} vanishes, so the M-step maximizes
$\E_{q^{(t)}}[\log f_\theta(x, U)]$, which separates over $\theta$
because all summations move outside the logarithm. The Lagrangian derivation extends
to the chain-structured statistics $\gamma_t, \xi_t$, giving the
closed-form updates
\begin{equation}
  \hat\pi_k = \gamma_1(k), \qquad
  \hat A_{jk} = \frac{\sum_t \xi_t(j,k)}{\sum_t \gamma_t(j)}, \qquad
  \hat\mu_k = \frac{\sum_t \gamma_t(k)\,x_t}{\sum_t \gamma_t(k)}, \qquad
  \hat\sigma_k^2 = \frac{\sum_t \gamma_t(k)\,(x_t - \hat\mu_k)^2}{\sum_t \gamma_t(k)}.
  \label{eq:bw-mstep}
\end{equation}

\begin{algorithm}[ht]
  \caption{Base EM (Baum--Welch)}
  \begin{algorithmic}[1]
    \State Initialize $\theta^{(0)}$
    \Repeat
    \State \textbf{E-step:} Run forward--backward to compute
    $\gamma_t(k)$, $\xi_t(j,k)$ via Eqs.~\eqref{eq:gamma}--\eqref{eq:xi}
    \State \textbf{M-step:} Update $\theta^{(t)}$ via Eq.~\eqref{eq:bw-mstep}
    \Until{$|\log f_{\theta^{(t)}}(x) - \log f_{\theta^{(t-1)}}(x)| < \varepsilon$}
  \end{algorithmic}
\end{algorithm}

\subsection{Variational EM (VEM)}

VEM restricts $q$ to a tractable family $\mathcal{Q}$ to handle settings
where the exact posterior is intractable. Start by adopting the mean-field family
$q(z) = \prod_{t=1}^T q_t(z_t)$, which discards the chain dependencies
preserved by Eqs.~\eqref{eq:gamma}--\eqref{eq:xi}. The E-step is the
Coordinate Ascent Variational Inference (CAVI) sweep: cycling through $t = 1, \ldots, T$, update
\begin{equation}
  \log q_t(k) \;\propto\;
  \log \mathcal{N}(x_t; \mu_k, \sigma_k^2)
  \;+\; \sum_j q_{t-1}(j)\,\log A_{jk}
  \;+\; \sum_j q_{t+1}(j)\,\log A_{kj},
  \label{eq:vem-q}
\end{equation}
which is the HMM specialization of the generic CAVI rule
$\log q_j \propto \E_{-j}[\log f(\Theta_j, \Theta_{-j}, x)]$.
After CAVI converges, the chain-aware
$\gamma_t, \xi_t$ in Eq.~\eqref{eq:bw-mstep} are replaced by their
mean-field surrogates
\begin{equation}
  \tilde\gamma_t(k) = q_t(k), \qquad
  \tilde\xi_t(j,k)  = q_t(j)\, q_{t+1}(k),
  \label{eq:vem-stats}
\end{equation}
and the M-step then applies Eq.~\eqref{eq:bw-mstep} with these
substitutions. Because $q$ is restricted to $\mathcal{Q}$, the $\KL$
term in Eq.~\eqref{eq:elbo} does not vanish and VEM maximizes a strict
lower bound on $\log f_\theta(x)$.

\begin{algorithm}[ht]
  \caption{Variational EM (mean-field)}
  \begin{algorithmic}[1]
    \State Initialize $\theta^{(0)}$ and $q^{(0)}(z) = \prod_t q_t^{(0)}(z_t)$
    \Repeat
    \State \textbf{E-step:} Sweep CAVI updates Eq.~\eqref{eq:vem-q} until
    the ELBO converges; form $\tilde\gamma_t, \tilde\xi_t$ via
    Eq.~\eqref{eq:vem-stats}
    \State \textbf{M-step:} Update $\theta^{(t)}$ via
    Eq.~\eqref{eq:bw-mstep} with $\tilde\gamma_t, \tilde\xi_t$ in place
    of $\gamma_t, \xi_t$
    \Until{$|F(q^{(t)},\theta^{(t)}) - F(q^{(t-1)},\theta^{(t-1)})| < \varepsilon$}
  \end{algorithmic}
\end{algorithm}

\subsection{Expectation conditional maximization (ECM)}

ECM replaces the joint M-step of
Eq.~\eqref{eq:bw-mstep} with a sequence of conditional maximizations
over parameter blocks. We partition $\theta$ into a dynamics block
$\theta_d = (\pi, A)$ and a $(\mu, \sigma^2)$ block
$\theta_e = (\mu_{1:K}, \sigma_{1:K}^2)$. The first CM-step uses
$\gamma_t, \xi_t$ from Eqs.~\eqref{eq:gamma}--\eqref{eq:xi} computed
under $\theta^{(t-1)}$ and updates only the dynamics:
\begin{equation}
  \hat\pi_k = \gamma_1(k), \qquad
  \hat A_{jk} = \frac{\sum_t \xi_t(j,k)}{\sum_t \gamma_t(j)}.
  \label{eq:ecm-cm1}
\end{equation}
We then refresh the E-step under
$(\hat\pi, \hat A, \mu^{(t-1)}, \sigma^{(t-1)})$ to obtain new marginals
$\gamma_t^{\,\prime}$, and the second CM-step updates $(\mu, \sigma^2)$
using these refreshed posteriors:
\begin{equation}
  \hat\mu_k = \frac{\sum_t \gamma_t^{\,\prime}(k)\, x_t}{\sum_t \gamma_t^{\,\prime}(k)},
  \qquad
  \hat\sigma_k^2 = \frac{\sum_t \gamma_t^{\,\prime}(k)\,(x_t - \hat\mu_k)^2}{\sum_t \gamma_t^{\,\prime}(k)}.
  \label{eq:ecm-cm2}
\end{equation}
Each CM-step is a partial maximization of $F(q, \theta)$, so ECM
inherits the monotone ascent of EM; the cost is
one additional forward--backward pass per outer iteration.

\begin{algorithm}[ht]
  \caption{Expectation Conditional Maximization (ECM)}
  \begin{algorithmic}[1]
    \State Initialize $\theta^{(0)} = (\pi^{(0)}, A^{(0)}, \mu^{(0)}, \sigma^{(0)})$
    \Repeat
    \State \textbf{E-step 1:} Run forward--backward under
    $\theta^{(t-1)}$ to compute $\gamma_t, \xi_t$ via
    Eqs.~\eqref{eq:gamma}--\eqref{eq:xi}
    \State \textbf{CM-step 1:} Update $(\pi, A)$ via Eq.~\eqref{eq:ecm-cm1}
    \State \textbf{E-step 2:} Run forward--backward under
    $(\hat\pi, \hat A, \mu^{(t-1)}, \sigma^{(t-1)})$ to obtain
    $\gamma_t^{\,\prime}$
    \State \textbf{CM-step 2:} Update $(\mu, \sigma^2)$ via Eq.~\eqref{eq:ecm-cm2}
    \Until{$|\log f_{\theta^{(t)}}(x) - \log f_{\theta^{(t-1)}}(x)| < \varepsilon$}
  \end{algorithmic}
\end{algorithm}

\subsection{Generalized EM (GEM)}

GEM relaxes the M-step to require only a non-decreasing update
of $F(q^{(t)}, \theta)$. We instantiate this
with a damped step toward the Base EM maximizer
$\hat\theta^{(t)}$ defined by Eq.~\eqref{eq:bw-mstep}.
Here the update is
\begin{equation}
  \eta^{(t)} \;=\; (1 - \alpha)\, \eta^{(t-1)} \;+\; \alpha\,
  \hat\eta^{(t)},
  \qquad \alpha \in (0, 1],
  \label{eq:gem-update}
\end{equation}
And any $\alpha \in (0, 1]$ yields a non-decreasing update
and therefore a valid GEM step; $\alpha = 1$
recovers Base EM. We use $\alpha = 0.5$ throughout to expose the
convergence-rate cost of an inexact M-step.

\begin{algorithm}[ht]
  \caption{Generalized EM with damped M-step}
  \begin{algorithmic}[1]
    \Require Damping coefficient $\alpha \in (0, 1]$
    \State Initialize $\theta^{(0)}$
    \Repeat
    \State \textbf{E-step:} Run forward--backward to compute
    $\gamma_t, \xi_t$ via Eqs.~\eqref{eq:gamma}--\eqref{eq:xi}
    \State \textbf{M-step:} Compute $\hat\theta^{(t)}$
    via Eq.~\eqref{eq:bw-mstep}; update $\theta^{(t)}$ via
    Eq.~\eqref{eq:gem-update}
    \Until{$|\log f_{\theta^{(t)}}(x) - \log f_{\theta^{(t-1)}}(x)| < \varepsilon$}
  \end{algorithmic}
\end{algorithm}


\pagebreak

% ============================================================
\section{Experiments}
% ============================================================

\paragraph{Setup.} We start by generating some HMMs, then simulate a large number of observations from each truth and run the four
methods for from a shared initialization. Convergence
is reported as the excess log-likelihood
$\log f_\theta(x) - \log f_{\theta^\star}(x)$ so that trajectories
aggregate across the random truths.

\paragraph{Convergence and $\hat\mu$ recovery.}
Figure~\ref{fig:em-comparison} shows convergence per iteration. Base
EM, ECM, and VEM close the gap to the MLE within $10$
iterations or so while GEM ($\alpha = 0.5$) takes roughly twice as long. The $\hat\mu$ error in panel (b)
shows the log-likelihood trajectory for each method.

\begin{figure}[t]
  \centering
  \begin{tikzpicture}
    \begin{groupplot}[
        group style={group size=2 by 1, horizontal sep=1.6cm},
        emaxis,
      ]
      \nextgroupplot[
        title={(a) Convergence vs.\ iteration},
        xlabel={EM iteration},
        ylabel={$\log f_\theta(x) - \log f_{\theta^\star}(x)$},
        xmin=0, xmax=30,
        legend style={at={(0.97,0.04)}, anchor=south east},
      ]
      \embandline{../code/convergence_iter.csv}{baseColor}{base}{iter}{Base EM}
      \embandline{../code/convergence_iter.csv}{vemColor}{vem}{iter}{VEM}
      \embandline{../code/convergence_iter.csv}{ecmColor}{ecm}{iter}{ECM}
      \embandline{../code/convergence_iter.csv}{gemColor}{gem}{iter}{GEM}

      \nextgroupplot[
        title={(b) $\hat\mu$ recovery},
        xlabel={EM iteration},
        ylabel={$\tfrac{1}{K}\sum_k |\hat\mu_k - \mu_k^\star|$},
        ymode=log, xmin=0, xmax=30,
        legend style={at={(0.97,0.96)}, anchor=north east},
      ]
      \embandline{../code/mu_error.csv}{baseColor}{base}{iter}{Base EM}
      \embandline{../code/mu_error.csv}{vemColor}{vem}{iter}{VEM}
      \embandline{../code/mu_error.csv}{ecmColor}{ecm}{iter}{ECM}
      \embandline{../code/mu_error.csv}{gemColor}{gem}{iter}{GEM}
    \end{groupplot}
  \end{tikzpicture}
  \caption{EM variants on Gaussian HMMs with random ground truth.
    (a) Excess log-likelihood. (b) $\hat\mu$ error.}
  \label{fig:em-comparison}
\end{figure}

\paragraph{Distance to the true distribution.}
Figure~\ref{fig:em-dist} plots a composite distance between each
recovered HMM and the ground truth.

\begin{figure}[t]
  \centering
  \begin{tikzpicture}
    \begin{axis}[
        emaxis,
        width=0.66\linewidth, height=0.42\linewidth,
        title={Distance to true HMM vs.\ iteration},
        xlabel={EM iteration},
        ylabel={Composite distance to $\theta^\star$},
        xmin=0, xmax=40,
        legend style={at={(0.97,0.96)}, anchor=north east},
      ]
      \embandline{../code/dist_to_truth.csv}{baseColor}{base}{iter}{Base EM}
      \embandline{../code/dist_to_truth.csv}{vemColor}{vem}{iter}{VEM}
      \embandline{../code/dist_to_truth.csv}{ecmColor}{ecm}{iter}{ECM}
      \embandline{../code/dist_to_truth.csv}{gemColor}{gem}{iter}{GEM}
    \end{axis}
  \end{tikzpicture}
  \caption{Distance to $\theta^\star$ across iterations.}
  \label{fig:em-dist}
\end{figure}

\paragraph{Recovered marginal density.}
Figure~\ref{fig:em-density} compares the stationary marginal
$p_\theta(x) = \sum_k \pi^\infty_k \, \mathcal{N}(x; \mu_k, \sigma_k^2)$
under each method to the truth on a single random selection, with the
empirical histogram of the observations as backdrop. Base
EM, ECM, and VEM lay close to the ground truth, while GEM does not quickly close the distance.

\begin{figure}[t]
  \centering
  \begin{tikzpicture}
    \begin{axis}[
        emaxis,
        width=0.78\linewidth, height=0.40\linewidth,
        title={Recovered marginal density vs.\ truth},
        xlabel={$x$}, ylabel={Density},
        ymin=0,
        legend style={at={(0.97,0.96)}, anchor=north east},
      ]
      \addplot[ybar interval, fill=gray!30, draw=gray!55,
        fill opacity=0.6, forget plot]
      table[x=x, y=density, col sep=comma] {../code/density_hist.csv};
      \addplot[black, very thick, dashed]
      table[x=x, y=truth_pdf, col sep=comma] {../code/density_curves.csv};
      \addlegendentry{Truth}
      \addplot[baseColor, thick]
      table[x=x, y=base_pdf, col sep=comma] {../code/density_curves.csv};
      \addlegendentry{Base EM}
      \addplot[vemColor, thick]
      table[x=x, y=vem_pdf, col sep=comma] {../code/density_curves.csv};
      \addlegendentry{VEM}
      \addplot[ecmColor, thick]
      table[x=x, y=ecm_pdf, col sep=comma] {../code/density_curves.csv};
      \addlegendentry{ECM}
      \addplot[gemColor, thick]
      table[x=x, y=gem_pdf, col sep=comma] {../code/density_curves.csv};
      \addlegendentry{GEM}
    \end{axis}
  \end{tikzpicture}
  \caption{Recovered marginal density.}
  \label{fig:em-density}
\end{figure}


\pagebreak
% ============================================================
\section{Discussion}
% ============================================================


Each of the four tested algorithms converges fairly quickly to the ground truth but not all at the same rate.
Base EM, ECM, and VEM all reach the MLE within about $10$ iterations, while GEM ($\alpha = 0.5$) takes roughly twice as long.
I found that it is preferable to use Base EM when
feasible, and to utilize a variant only algorithm when a
harder model forces such a choice.


% ============================================================
\bibliographystyle{plainnat}
\bibliography{references}
\end{document}
